\chapter{Risoluzione dei problemi (FAQ)}
\label{ch:faq}
\index{FAQ}
\index{Risoluzione dei problemi}

\epigraph{``In teoria non c'\`e differenza tra teoria e pratica.
Nella pratica, s\`i.''}{\textsc{attribuito a} Yogi Berra}

\lettrine[lines=3]{Q}{uesto} capitolo raccoglie in un unico posto
i problemi che capitano pi\`u spesso, con la soluzione a portata
di mano. \`E organizzato per \emph{sintomo}: cerca la frase che
somiglia a quello che vedi e segui i passi. Se un termine non ti
\`e chiaro, il \hyperref[ch:terminologia_didattica]{glossario dei
termini scolastici} lo spiega; per i dettagli operativi rimandiamo
alla \hyperref[ch:getting_started]{guida rapida} e alle
\hyperref[ch:workflow_tipici]{procedure tipiche}.

\begin{avvertenzabox}
Regola d'oro: quasi tutti i problemi si diagnosticano dal
\emph{pannello di log} (i messaggi che scorrono durante un
lancio) e dal banner di stato in alto. Un messaggio rosso non
va ignorato: dice quasi sempre \emph{cosa} \`e mancato.
\end{avvertenzabox}

\section{Avvio e installazione}

\begin{description}

\item[Il browser si apre ma vedo una pagina bianca.]
Apri la console del browser (\texttt{F12}): nella quasi totalit\`a
dei casi il backend non risponde. Verifica
\texttt{http://localhost:8000/api/health}: deve restituire un
piccolo JSON \texttt{\{"status":"ok"\}}. Se non risponde, il
processo del backend non \`e partito o si \`e chiuso.

\item[``Backend non raggiungibile''.]
Il dashboard contatta \texttt{http://localhost:8000/api/health};
se non risponde compare un banner rosso. Controlla che il
processo del server (\texttt{uvicorn}) sia ancora attivo e che
la porta \texttt{8000} non sia occupata da un altro programma.

\item[``La porta 5173 (o 8000) \`e gi\`a in uso''.]
Un'istanza precedente non \`e stata chiusa. Su Linux/macOS
\texttt{./webui/stop.sh} ferma tutto in modo pulito; su Windows
chiudi le finestre lasciate aperte o termina il processo
\texttt{node.exe}/\texttt{python.exe} dal Task Manager. Poi
riavvia con \texttt{start.sh}/\texttt{start.bat}.

\item[Il primo avvio \`e lentissimo.]
\`E normale: al primo lancio vengono create l'ambiente Python e
le dipendenze del frontend (qualche minuto). I lanci successivi
sono immediati. Vedi il capitolo \hyperref[ch:launcher]{Avviare e
fermare il programma}.

\end{description}

\section{Dati, import ed export}

\begin{description}

\item[Il dashboard \`e vuoto (tutti zeri).]
Non c'\`e ancora una scuola in memoria. Hai due strade: inserire
i tuoi dati (Ore \textrightarrow{} Materie \textrightarrow{}
Docenti \textrightarrow{} Classi \textrightarrow{} Aule
\textrightarrow{} Cattedre), oppure importare un dataset di
esempio dal dashboard per fare una prova. Le procedure sono nel
capitolo \hyperref[ch:workflow_tipici]{Procedure tipiche}.

\item[``Import error'' / il file non viene accettato.]
Il file non rispetta il formato atteso (colonne o fogli
mancanti). Usa i \emph{modelli} scaricabili dalla pagina di
import: i fogli \emph{Istruzioni}, \emph{Esempi} e \emph{Dati}
mostrano esattamente le colonne. Prima di scrivere davvero, lancia
l'import in modalit\`a \emph{anteprima (dry-run)}: mostra riga per
riga cosa verrebbe creato senza toccare il database.

\begin{avvertenzabox}
Un import ``a sostituzione'' (\emph{replace}) di un intero
database \textbf{cancella} i dati esistenti. Fai prima un
backup/esporta, e in caso di dubbio usa la modalit\`a
\emph{anteprima}.
\end{avvertenzabox}

\item[Ho rinominato un docente/una classe e l'orario sembra rotto.]
Le lezioni gi\`a piazzate fanno riferimento ai nomi: dopo una
rinomina possono restare ``orfane''. Rilancia l'ottimizzazione,
oppure ripristina il nome precedente.

\end{description}

\section{Vincoli}

\begin{description}

\item[Ho cambiato la matrice di indisponibilit\`a ma non vedo effetto.]
Le modifiche ai vincoli entrano in gioco al \emph{successivo}
lancio della pipeline: non spostano lezioni gi\`a piazzate.
Rilancia dalla pagina di ottimizzazione (\emph{Workflow}).

\item[``DSL: parser error''.]
Se scrivi vincoli nel linguaggio DSL, le parole chiave sono
\emph{minuscole} e \emph{case-sensitive}: \texttt{forall},
\texttt{where}, \texttt{in}, \texttt{exists}, \texttt{count}. Gli
identificatori (nomi di docenti, materie, classi) sono valori
letterali, senza virgolette. Il capitolo
\hyperref[part:dsl]{Il DSL dei vincoli} ha la grammatica completa.

\item[Il modello risulta ``non fattibile'' (infeasible).]
Vuol dire che i vincoli \emph{HARD} si contraddicono a vicenda:
non esiste alcun orario che li rispetti tutti. Non \`e un errore
del programma, \`e la scuola che chiede l'impossibile. Usa il
\emph{Feasibility Check} nella pagina Vincoli per individuare il
sottoinsieme minimo di vincoli in conflitto, poi ammorbidiscine
uno (da \emph{HARD} a \emph{SOFT}) o rimuovilo. Vedi il capitolo
\hyperref[ch:vincoli]{Vincoli}.

\end{description}

\section{Ottimizzazione e run}

\begin{description}

\item[La pipeline si blocca su ``Phase A''.]
Su una scuola piccola Phase A finisce in pochi secondi. Se supera
la mezza minuto, spesso il database \`e stato importato male
(cattedre senza ore assegnate): reimporta i dati e riprova. Su
scuole molto grandi \`e invece normale che serva pi\`u tempo:
distingui i due casi dal messaggio (``nessuna soluzione'' \`e
diverso da ``non risolto nel tempo'').

\item[``Pipeline completata ma orario vuoto''.]
Controlla sul dashboard che il numero di lezioni sia positivo. Se
\`e zero, la pipeline \`e probabilmente fallita: leggi i messaggi
rossi nel pannello di log (fase Phase B). Con il \emph{controllo
di copertura} attivo, un orario incompleto viene ora segnalato
come fallimento invece di essere salvato in silenzio.

\item[Come interrompo un lancio partito per errore?]
Ogni run pu\`o essere annullato: nella pagina \emph{Runs} il
pulsante di annullamento ferma il lavoro. L'annullamento \`e
\emph{cooperativo}, quindi pu\`o richiedere qualche secondo per
rientrare da un passaggio lungo.

\item[Ho lanciato due ottimizzazioni e la seconda resta ``in coda''.]
\`E voluto: per non saturare la macchina il sistema esegue un
lancio pesante alla volta (\emph{admission control}); i lanci
eccedenti restano \emph{pending} e partono appena si libera lo
slot.

\item[Un run \`e rimasto ``running/pending'' per sempre dopo un riavvio.]
Se il backend si \`e chiuso a met\`a, il run resta ``appeso''. Al
riavvio il sistema fa pulizia automatica e lo marca come
\emph{failed}: non \`e un dato reale ma un residuo, puoi
eliminarlo.

\end{description}

\section{Orario e modifiche manuali}

\begin{description}

\item[Come sposto una lezione a mano?]
Nella pagina \emph{Orario} trascina la lezione su uno slot vuoto;
un click sulla lezione apre le azioni (sposta, elimina, svincola).
Le lezioni ``svincolate'' finiscono nel pool laterale, da cui puoi
ritrascinarle sul calendario.

\item[Vedo un avviso ``lezioni fuori dalla griglia''.]
Alcune lezioni sono su un'ora che non \`e pi\`u configurata nel
tab \emph{Ore} (magari l'hai tolta dopo aver generato l'orario):
non hanno una cella dove comparire. Riabilita l'ora corrispondente
oppure sposta/elimina quelle lezioni.

\item[Una modifica manuale viola un vincolo?]
Il calendario mostra in anteprima i conflitti \emph{soft} mentre
trascini; le celle bloccate (\emph{lock}) non si possono
sovrascrivere. Un lock \`e un vincolo \emph{HARD}: il solver lo
rispetter\`a in ogni rilancio.

\end{description}

\section{Prestazioni}

\begin{description}

\item[Tutto \`e lento e la macchina \`e satura.]
Un altro programma sta consumando la CPU, oppure stai lanciando
un profilo molto grande (\texttt{huge}/\texttt{superhuge}). Per un
uso reale una scuola sta comodamente nei profili medi; i profili
enormi servono ai benchmark. Chiudi gli altri processi pesanti
prima di un lancio importante.

\item[``Database is locked''.]
Capita se molte scritture concorrono sullo stesso database SQLite.
Il sistema attende gi\`a fino a qualche secondo prima di
arrendersi; se persiste, evita di lanciare un'ottimizzazione
mentre importi o svuoti il database. Per usi con pi\`u utenti in
contemporanea valuta PostgreSQL.

\item[I risultati cambiano tra un lancio e l'altro a parit\`a di dati.]
Il solver \`e reso \emph{riproducibile} da un seme fisso
(\texttt{PITANTUM\_SOLVER\_SEED}); se vuoi risultati
bit-identici forza anche la modalit\`a deterministica
(\texttt{PITANTUM\_DETERMINISTIC}), che usa un solo worker.

\end{description}

\begin{biblioparvabox}
Non hai trovato il tuo problema? Il pannello di log di un run e la
console del browser (\texttt{F12}) sono la prima fonte di
informazione. In seconda battuta, apri una \emph{issue} sul
repository (vedi l'appendice \hyperref[ch:repository_github]{Il
repository GitHub}).
\end{biblioparvabox}
